\documentclass[10pt]{article}
\usepackage[a4paper,margin=26mm]{geometry}
\usepackage{amsmath,amssymb,amsthm}
\usepackage[T1]{fontenc}
\usepackage{lmodern,microtype}
\usepackage[colorlinks=true,allcolors=blue]{hyperref}
\hypersetup{pdftitle={Nonmonotone Argument and Radial Angle of a Support Slit},
 pdfauthor={Alec Kriebel},pdfsubject={Unrefereed AI-assisted research note}}
\setlength{\emergencystretch}{2em}
\newtheorem{theorem}{Theorem}
\newtheorem{lemma}[theorem]{Lemma}
\newcommand{\eps}{\varepsilon}
\newcommand{\D}{\mathbb D}
\newcommand{\C}{\mathbb C}
\newcommand{\Rea}{\operatorname{Re}}
\newcommand{\Ima}{\operatorname{Im}}
\title{Nonmonotone Argument and Radial Angle of a Support Slit}
\author{Alec Kriebel\\
 \small ORCID: \href{https://orcid.org/0009-0001-9320-500X}{0009-0001-9320-500X}}
\date{13 October 2026}
\begin{document}
\maketitle
\begin{abstract}
We give negative answers to the two slit-monotonicity questions of
P.~L.~Duren recorded as Problem~6.72 in Hayman--Lingham.
Every global maximizer of the real part of
\(a_2+i10^{-24}(a_4-3a_3)\) in the normalized schlicht class
omits an arc with nonmonotone continuous argument and nonmonotone radial
angle, under both signed and ordinary unsigned conventions.
The proof uses classical support-point geometry, a quantitative area-theorem
estimate for the actual slit tip, and controlled inversion of the
Schiffer differential. It establishes actual extremizers by compactness
and does not assume that an arbitrary trajectory is an extremal slit.
No formula for the maximizing function or uniqueness is required.
This is an unrefereed AI-assisted research note.
\end{abstract}

\section{The question and precise statement}
Let \(S\) consist of analytic injective maps on \(\D=\{|z|<1\}\)
normalized by
\[
                 f(z)=z+\sum_{n\ge2}a_n(f)z^n.
\]
A support point maximizes \(\Rea L\) for a continuous complex-linear
functional on \(H(\D)\) whose real part is nonconstant on \(S\).
Duren's questions ask whether the omitted analytic arc of a support
point must have monotone argument and monotone angle between its radius
and tangent vectors \cite[Problem 6.72]{HL}. Monotone here allows either
direction. Argument means a continuous real-valued lift; tangent vectors
are directed outward toward infinity.

\begin{theorem}\label{main}
Set
\[
 \eps=10^{-24},\qquad \eta=10^{-12},\qquad
 L(f)=a_2(f)+i\eps\bigl(a_4(f)-3a_3(f)\bigr).
\]
Global maximizers of \(\Rea L\) on \(S\) exist. Every such maximizer
omits an arc \(\Gamma\) on which continuous argument, signed outward
radial angle, and the ordinary nonnegative outward radial angle
are all nonmonotone.
\end{theorem}

We use compactness of \(S\), the coefficient inequalities
\(|a_n|\le n\) \cite{deBranges}, and the exterior area theorem.
The classical support-slit theorem gives a single analytic omitted arc
from a finite tip to infinity, with increasing modulus, satisfying
\begin{equation}\label{schiffer}
 Q(w)\,dw^2>0,\qquad
 Q(w)=\frac{\Phi(w)}{w^2},\qquad
 \Phi(w)=L\left(\frac{f^2}{f-w}\right).
\end{equation}
This is Duren's convention \cite[pp.352--355]{Duren}.
Pfluger's original geometric proof is in
\cite[\S\S4.1--4.5]{Pfluger}; his differential has the opposite sign.
We use these established theorems with their stated hypotheses, and
derive the needed local trajectory uniqueness below.

\section{Uniform control of actual extremizers and their tips}
Coefficient maps are continuous on \(H(\D)\) by Cauchy's formula.
Compactness supplies a maximizer \(f\). The Koebe function
\(k(z)=z/(1-z)^2\) and the identity give real functional values \(2\)
and \(0\), so every maximizer is a genuine support point.
Coefficient bounds imply
\begin{equation}\label{a2}
 \Rea a_2\ge2-13\eps,\qquad
 |a_2-2|^2\le52\eps,\qquad
 0\le1-|a_2|^2/4\le13\eps.
\end{equation}
In particular \(|a_2-2|<8\eta\).

The normalized odd root \(g(z)^2=f(z^2)\) is univalent.
Indeed equality of two \(g\)-values implies equality of the squared
inputs; oddness excludes the opposite pair except at zero.
The exterior map
\[
 \frac1{g(1/\zeta)}
 =\zeta-\frac{a_2}{2}\zeta^{-1}
       +\left(\frac{3a_2^2}{8}-\frac{a_3}{2}\right)\zeta^{-3}+\cdots
\]
therefore satisfies the area theorem. It gives
\[
 |a_3-3a_2^2/4|\le2\sqrt{13/3}\,\eta<5\eta,\qquad
 |a_3-3|<29\eta.
\]
With
\[
 A=2a_3+a_2^2-6a_2,\qquad B=3a_2-3,
\]
we obtain
\begin{equation}\label{AB}
             |A+2|<106\eta,\qquad |B-3|<24\eta.
\end{equation}
The first follows from
\(A+2=2(a_3-3)+(a_2-2)(a_2-4)\) and \(|a_2-4|\le6\).
All bounds hold uniformly for every maximizer.

We next bound the \emph{actual} tip, without assuming continuity of
tips under perturbation. Write
\[
 H(z)=\sqrt{z/f(z)}=1-\frac{a_2}{2}z+\sum_{n\ge2}c_nz^n,
 \qquad H_0(z)=1-z.
\]
The same exterior map equals \(\zeta H(\zeta^{-2})\), so
\[
                  \sum_{n\ge2}(2n-1)|c_n|^2\le13\eps.
\]
For \(|z|\le r=2/3\), weighted Cauchy--Schwarz gives
\begin{equation}\label{H}
                 |H-H_0|<5\eta,\qquad |H'-H_0'|<10\eta.
\end{equation}
The two tail sums are bounded respectively by
\[
 \sum_{n\ge2}\frac{r^{2n}}{2n-1}
 \le\frac{r^4}{3(1-r^2)}=\frac{16}{135}<\frac14,
 \quad
 \sum_{n\ge2}\frac{n^2r^{2n-2}}{2n-1}
 \le(1-r^2)^{-2}-1=\frac{56}{25}<\frac94 .
\]
Together with \(\sqrt{13}<4\), these prove \eqref{H}.
At \(z_*=-2/3\), \(H_0=5/3\) and
\[
 \frac32<|H|<\frac{11}{6},\qquad
 |H+H_0|\le\frac{10}{3}+5\eta<\frac72.
\]
The identities \(f=z/H^2\) and \(f'=(H-2zH')/H^3\) give
\begin{equation}\label{values}
             |f(z_*)-k(z_*)|<2\eta,\qquad
             |f'(z_*)-k'(z_*)|<8\eta .
\end{equation}
Explicitly, the reciprocal-square error is at most
\((14/5)\eta\), before multiplication by \(|z_*|=2/3\).
The derivative numerator differs from \(1/3\) by less than \(19\eta\),
contributing less than \(6\eta\) after division by \(|H|^3\).
The remaining contribution is less than \(2\eta\), since
\[
 |H^2+HH_0+H_0^2|<11,\qquad |H|^3H_0^3>125/8.
\]

For a conformal map \(f:\D\to\Omega\), Schwarz's lemma applied to the
inverse on any inscribed disk about \(f(z)\) gives
\[
          \operatorname{dist}(f(z),\partial\Omega)
          \le(1-|z|^2)|f'(z)|.
\]
The direction of this inequality is essential here.
Using \(k(z_*)=-6/25,\ k'(z_*)=9/125\), \eqref{values} gives
\begin{equation}\label{tip}
 r_{\rm tip}=\min_{w\in\Gamma}|w|
 \le |f(z_*)|+\frac59|f'(z_*)|
 <\frac7{25}+7\eta<\frac3{10}.
\end{equation}
The omitted arc is closed, and its increasing modulus makes its
minimum the finite tip.

\section{The natural coordinate and the actual omitted subarc}
Coefficient extraction through degree four in
\(-f^2/w-f^3/w^2-f^4/w^3+\cdots\) gives exactly
\begin{equation}\label{Q}
 Q(w)=-\frac a{w^3}-\frac b{w^4}-\frac c{w^5},
 \quad a=1+i\eps A,\quad b=i\eps B,\quad c=i\eps .
\end{equation}
For omitted \(w\), this is \eqref{schiffer}; elsewhere it is its
meromorphic continuation. In particular \(a=L(f^2)\ne0\).
Under the double cover \(w=-v^{-2}\),
\[
 Q(w)\,dw^2=4P(v)\,dv^2,\qquad P(v)=a-bv^2+cv^4.
\]
Let \(P_0(v)=-2-3v^2+v^4\). From \eqref{AB}, on \(|v|\le3\),
\begin{equation}\label{P}
 |P-1-i\eps P_0|\le322\eps\eta,\quad
 |P_0|\le110,\quad |P-1|\le111\eps .
\end{equation}
Choose the analytic square root \(p=\sqrt P\) with positive real part
and put \(G(v)=\int_0^v p(s)\,ds\).
The identity \((p-1)(p+1)=P-1\) gives
\[
                 |p-1|\le111\eps,\qquad |G(v)-v|\le333\eps.
\]
For every \(|t|<2\), Rouché on \(|v|=3\) gives exactly one root
of \(G(v)=t\). Since \(G'=p\ne0\), the roots patch to a holomorphic
inverse \(V\), with
\begin{equation}\label{V}
 V(0)=0,\quad V'(0)=a^{-1/2},\quad
 |V(t)-t|\le333\eps,\quad V'(t)=1/p(V(t)).
\end{equation}
Both \(G\) and \(V\) are odd.

For \(0<t\le T=7/4\), define \(W(t)=-V(t)^{-2}\). It is regular,
\begin{equation}\label{W}
 Q(W)W'^2=4,\qquad W(t)\sim-a/t^2,\qquad
 |W(t)|>25/81>3/10,
\end{equation}
because \(|V(t)|<9/5\).
It has no self-intersection: equality of two \(W\)-values implies
\(V(t_1)=\pm V(t_2)\), and the odd inverse relation forces
\(t_1=\pm t_2\).

We now identify \(W\) with the actual omitted slit.
A distant tail of \(\Gamma\) lifts to \(v\to0\) in the double cover.
The positive trajectory condition says \(\Im G\) is constant; its
limiting value is zero. The two real \(G\)-rays project to the same
geometric \(w\)-arc by oddness of \(V\). Thus \(W\) and \(\Gamma\)
coincide near infinity.
The set
\[
             E=\{t\in(0,T]:W(t)\in\Gamma\}
\]
is nonempty and relatively closed. At every member, \eqref{tip}
and \eqref{W} exclude the actual tip, with radius gap \(7/810\).
Moreover \(Q\) is finite and nonzero there. A local square-root
primitive of \(Q\) makes positive trajectories horizontal lines,
so local uniqueness makes \(E\) relatively open. Connectedness
gives \(E=(0,T]\). No converse extremality criterion for arbitrary
trajectories has been used.

Put \(Z=tV'/V\). Independently of the fine expansion below,
\begin{equation}\label{orientation}
 Z(t)=\frac{G(V(t))/V(t)}{p(V(t))},\qquad
 |Z(t)-1|\le\frac{222\eps}{1-111\eps}<223\eps
                    \quad(0<t\le T).
\end{equation}
Here \(G(v)/v\) is the average of \(p\) on a radial segment.
Thus \(\Rea Z>0\), and
\(d\log|W|/dt=-2\Rea Z/t<0\): increasing \(t\) is inward,
so \(-W'\) is the outward tangent.

\section{Certified reversal of both angles}
Define
\[
 d(t)=t+\tfrac12t^3-\tfrac1{10}t^5,\qquad
 q(t)=d'(t)-d(t)/t=t^2-\tfrac25t^4 .
\]
We claim, uniformly for \(1\le t\le T\),
\begin{equation}\label{errors}
 |V-t-i\eps d(t)|\le487\eps\eta,\qquad
 |V'-1-i\eps d'(t)|\le162\eps\eta .
\end{equation}
Here are explicit analytic remainder bounds.
On the selected square-root branch,
\[
 \sqrt{1+x}-1-x/2=-\frac{x^2}{2(\sqrt{1+x}+1)^2}.
\]
Together with \eqref{P}, it gives
\[
 |p-1-(i\eps/2)P_0|\le162\eps\eta.
\]
Indeed the costs are \(161\eps\eta+(111^2/2)\eps^2\).
Since \(P_0=-2d'\), integration gives
\(G(v)=v-i\eps d(v)+R(v)\) with \(|R|\le486\eps\eta\).
The bound \(|d'|\le55\) on \(|v|\le3\) and \eqref{V}
make inversion cost at most \(18315\eps^2\), proving the first
inequality in \eqref{errors}.

Complex Taylor's integral remainder gives
\[
 |(1+x)^{-1/2}-1+x/2|\le2|x|^2\quad(|x|\le1/4).
\]
For example its second derivative has modulus at most
\((3/4)(4/3)^{5/2}<4\).
The coefficient uncertainty costs \(161\eps\eta\); the quadratic
term costs \(24642\eps^2\). Since \(|P_0'|\le126\), replacing
\(P_0(V)\) by \(P_0(t)\) costs \(20979\eps^2\).
The total is at most \(162\eps\eta\), proving the second inequality.
All absorptions are justified by
\[
 (111^2/2)\eta<1,\quad18315\eta<1,\quad45621\eta<1.
\]

Writing the two remainders in \eqref{errors} as \(R_V,R_{V'}\),
subtraction before division gives the exact identity
\[
 V(Z-1-i\eps q)=tR_{V'}-R_V-i\eps q(V-t).
\]
On this interval \(|V|>1/2\) and \(|q|<7\). Hence
\begin{equation}\label{Zerror}
 |Z-1-i\eps q|
 \le1541\eps\eta+4662\eps^2<2000\eps\eta .
\end{equation}
The exact values
\[
 q(1)=3/5,\qquad q(T)=-441/640,\qquad2000\eta=2\cdot10^{-9}
\]
therefore imply
\[
                  \Ima Z(1)>\eps/2,\qquad
                  \Ima Z(T)<-\eps/2 .
\]

For a continuous argument \(\theta\) of \(W\),
\[
       \theta'(t)=\Ima(W'/W)=-2\Ima Z(t)/t.
\]
Its derivative is negative at \(1\) and positive at \(T\), both
regular points of the actual slit. It cannot be monotone.
The signed outward radial angle is
\[
              \beta(t)=\arg((-W')/W)=\arg Z(t).
\]
The branch is fixed by \eqref{orientation}; it tends to zero as
\(t\downarrow0\), since \(V(t)=a^{-1/2}t+O(t^3)\).
It is positive at \(1\) and negative at \(T\).
A nondecreasing function with initial limit zero cannot take
the negative value, and a nonincreasing one cannot take the positive
value. Thus the signed angle is nonmonotone.
Continuity gives a zero between these two points, whereas
\(|\beta|\) is positive at both endpoints. The ordinary unsigned
angle is therefore nonmonotone even on that compact subarc.
An inward tangent gives its supplement, and the acute tangent-line
interpretation gives \(|\beta|\); neither restores monotonicity.
This proves Theorem~\ref{main}.

\section{Attribution, scope and verification}
The support-slit geometry and Schiffer method are classical.
Duren \cite{Duren} and Brown \cite{Brown} established monotonicity
for important special functionals; these results do not impose
monotonicity for the mixed fourth-degree functional used here.
Our argument controls every maximizer without a numerical extremal
solver, a unique maximizing selection, or boundary-tip continuity.
At \(\eps=0\) the sign reversal disappears.

A change of tangent direction concerns \(\arg w'\), whereas the
radius argument is \(\arg w\) and the radial angle is
\(\arg w'-\arg w\). Zemyan's earlier tangent-waver construction
\cite{Zemyan03} is consequently related prior work that must not be
confused with either target conclusion solely from its abstract.
His later abstract reports finitely many changes of direction
\cite{Zemyan07}. Neither published abstract is a substitute for
reading the complete construction or proof.

The primary-source priority audit is bounded. Complete texts of
several closely relevant leads remain unread at this draft stage,
including the Zemyan papers and the Hibschweiler and Ma--Zhang
papers identified in the accompanying audit. No historical-first
priority or still-open-in-2026 assertion is made. This paragraph
will be reconciled with the completed priority assessment before
the package is frozen for its fresh publication reviews.

The supporting archive will contain the manuscript source, immutable
authored input, independent analytic and quantitative derivations,
exact rational checks, source and priority receipts, and integrity
manifests. Finite checks supplement the uniform analytic proofs;
they are not formal certification of the classical theorems.
AI tools were used extensively in development, writing, reproduction
and adversarial verification. This preprint is unrefereed and has
received no human peer review. Internal reviews are AI-assisted.

\begin{thebibliography}{9}
\raggedright
\bibitem{HL}
W.~K.~Hayman and E.~F.~Lingham, \emph{Research Problems in Function
Theory (New Edition)}, arXiv:1809.07200v2 (2018), Problem~6.72, p.143.
\href{https://arxiv.org/abs/1809.07200v2}{arXiv:1809.07200v2}.
\bibitem{Duren}
P.~L.~Duren, Arcs omitted by support points of univalent functions,
\emph{Comment. Math. Helv.} \textbf{56} (1981), 352--365.
\href{https://doi.org/10.1007/BF02566218}{doi:10.1007/BF02566218}.
\bibitem{Pfluger}
A.~Pfluger, Lineare Extremalprobleme bei schlichten Funktionen,
\emph{Ann. Acad. Sci. Fenn. Ser. A I} \textbf{489} (1971), 1--32.
\href{https://doi.org/10.5186/aasfm.1971.489}{doi:10.5186/aasfm.1971.489}.
\bibitem{deBranges}
L.~de~Branges, A proof of the Bieberbach conjecture,
\emph{Acta Math.} \textbf{154} (1985), 137--152.
\href{https://doi.org/10.1007/BF02392821}{doi:10.1007/BF02392821}.
\bibitem{Brown}
J.~E.~Brown, A method for investigating geometric properties of
support points and applications,
\emph{Trans. Amer. Math. Soc.} \textbf{287} (1985), 285--291.
\href{https://doi.org/10.1090/S0002-9947-1985-0766220-8}{doi:10.1090/S0002-9947-1985-0766220-8}.
\bibitem{Zemyan03}
S.~M.~Zemyan, The existence of a waver point on the omitted arc of
a support point of the class \(S\),
\emph{Complex Variables Theory Appl.} \textbf{48} (2003), 791--796.
\href{https://doi.org/10.1080/02781070310001599359}{doi:10.1080/02781070310001599359}.
\bibitem{Zemyan07}
S.~M.~Zemyan, On the behavior of the omitted arc of a support point
of the class \(S\),
\emph{Complex Variables Elliptic Equations} \textbf{52} (2007), 175--181.
\href{https://doi.org/10.1080/17476930601065169}{doi:10.1080/17476930601065169}.
\end{thebibliography}
\smallskip\noindent\small
Version 1.0 draft. Author: Alec Kriebel. License: CC BY 4.0.
\end{document}
